%% TRANSCRIBED (PDF 121-145 = printed III-30..III-53)
%% NOTE ON THE SCAN: printed pages III-37 and III-38 are ABSENT from this
%% PDF (its PDF pages 128-129 contain duplicate copies of IV-37/IV-38,
%% which also occur at their proper place, PDF 182-183).  Those two
%% printed pages were transcribed from the complete scan of the same book
%% at working/serre-mcgill.pdf (its PDF pages 129-130; page images saved
%% as working/mcgill_III37.png and working/mcgill_III38.png) and verified
%% visually like every other page.
%% PDF 145 is the publisher's blank page (the appendix ends on III-53).

\apphead{Hodge--Tate decompositions and locally algebraic representations}%
  {HODGE--TATE DECOMPOSITIONS AND LOCALLY ALGEBRAIC REPRESENTATIONS}

Let $K$ be a field of characteristic zero, complete with respect to a
discrete valuation and with perfect residue field $k$ of characteristic
$p>0$.  In this Appendix we deal with Hodge--Tate decomposition of
$p$-adic abelian representations of $K$.

Sections \ref{sec:III.A.1} and \ref{sec:III.A.2} give invariance
properties of these decompositions under ground field extensions.
Special characters of $\Gal(\bar K/K)$ are defined in \ref{sec:III.A.4};
they are closely connected both with Hodge--Tate modules
(\ref{sec:III.A.4} and \ref{sec:III.A.5}) and local algebraicity
(\ref{sec:III.A.6}).  The proof of Tate's theorem
(cf.~\S\ref{sec:III.1.2}) is given in the last section.

\subsechead{A.1}{Invariance of Hodge--Tate decompositions}

Let $C$ be the completion of $\bar K$ (cf.~\S\ref{sec:III.1.2}); the
group $\Gal(\bar K/K)$ acts continuously on $C$.  Let $\chi$ be the
character of $\Gal(\bar K/K)$ into the group of $p$-adic units defined
in chap.~\ref{chap:I}, \S\ref{sec:I.1.2}.  Let $K'/K$ be a subextension
of $\bar K/K$ on which the valuation $\bar v$ of $\bar K$ is discrete;
this means that $K'$ is a finite extension of an unramified one of $K$.
Let $\hat K'$ denote the closure of $K'$ in $C$.

Let now $W$ be a finite dimensional $C$-vector space on which
$\Gal(\bar K/K)$ acts continuously and semi-linearly (see
\S\ref{sec:III.1.2}).  As before, we denote by $W^n$ (resp.\
$W^n_{K'}$) the $K$- (resp.\ $\hat K'$-) vector space defined by
\[
W^n=\{w\in W \mid s(w)=\chi(s)^n w \quad\text{for all}\quad
s\in\Gal(\bar K/K)\}
\]
(resp.
\[
W^n_{K'}=\{w\in W \mid s(w)=\chi(s)^n w \quad\text{for all}\quad
s\in\Gal(\bar K/K')\}).
\]
Let $W(n)=C\tens_K W^n$ and $W(n)'=C\tens_{\hat K'} W^n_{K'}$.
Identifying the modules $W(n)$ and $W(n)'$ with their canonical images
in $W$, we prove

\begin{theorem}[1]\label{thm:III.A.1.1}
The canonical map $\hat K'\tens_K W^n \to W^n_{K'}$ is a
$\hat K'$-isomorphism.
\end{theorem}

\begin{corollary}[1]\label{cor:III.A.1.1}
The Galois modules $W(n)$ and $W(n)'$ are equal.
\end{corollary}

Indeed, Theorem~\ref{thm:III.A.1.1} shows that $W^n$ and $W^n_{K'}$
generate the same $C$-vector subspace of $W$.

\begin{corollary}[2]\label{cor:III.A.1.2}
The Galois module $W$ is of Hodge--Tate type over $K$ if and only if it
is so over $\hat K'$.
\end{corollary}

\begin{proof}[Proof of Theorem~\ref{thm:III.A.1.1}.]
Note first that replacing the action of $\Gal(\bar K/K)$ on $W$ by
$(s,w)\mapsto \chi(s)^{-i}sw$, $i\in\Z$, just changes $W^n$ to
$W^{n+i}$.  This shifting process reduces the problem to the case
$n=0$; in that case, $W^n$ (resp.\ $W^n_{K'}$) is the set of elements of
$W$ which are invariant under $\Gal(\bar K/K)$ (resp.\ under
$\Gal(\bar K/K')$).  Note also that the injectivity of
$\hat K'\tens W^0 \to W^0_{K'}$ is trivial, since we know that
$C\tens_K W^0 \to W$ is injective (cf.~\S\ref{sec:III.1.2}).

On the other hand, an easy up-and-down argument shows that one can
assume $K'/K$ to be either \emph{finite Galois} or \emph{unramified
Galois}.  In both cases, since $\Gal(\bar K/K')$ acts trivially on
$W^0_{K'}$, we have a semi-linear action of $\Gal(K'/K)$ on
$W^0_{K'}$.  When $K'/K$ is finite, it is well known that this implies
that $W^0_{K'}$ is generated by the elements invariant by
$\Gal(K'/K)$, i.e.\ by $W^0$ (this is a non-commutative analogue of
Hilbert's ``Theorem 90'', cf.\ for instance \cite{serre-corps},
p.~159).

Let now $K'/K$ be unramified Galois and let $G$ be its Galois group.
Let $\hat O'$ denote the ring of integers of $\hat K'$.  Let $\Lambda$
be an $\hat O'$-lattice of $W^0_{K'}$ (i.e.\ a free $\hat O'$-submodule
of $W^0_{K'}$ of the same rank as $W^0_{K'}$).  Since $G$ acts
continuously on $W^0_{K'}$, the stabilizer in $G$ of $\Lambda$ is open,
hence of finite index, and the lattice $\Lambda$ has finitely many
transforms.  The sum $\Lambda^0$ of these transforms is invariant by
$G$.  Let $e_1,\ldots,e_N$ be a basis of $\Lambda^0$.  Let $s\in G$.
Then
\[
s(e_j)=\sum_{i=1}^{N} a_{ij}(s)e_i,\qquad a_{ij}\in\hat O',
\]
and the matrix $a(s)=(a_{ij}(s))$ belongs to $\GL(N,\hat O')$.  We have
$a(st)=a(s)s(a(t))$; this means that $a$ is a continuous $1$-cocycle on
$G$ with values in $\GL(N,\hat O')$.  Recall that two such cocycles
$a$ and $a'$ are said to be cohomologous if there exists
$b\in\GL(N,\hat O')$ such that $a'(s)=b^{-1}a(s)s(b)$ for all $s\in G$.
This is an equivalence relation on the set of cocycles and the
corresponding quotient space is denoted by $H^1(G,\GL(N,\hat O'))$.
In fact:

\begin{lemma}
$H^1(G,\GL(N,\hat O'))=\{1\}$.
\end{lemma}

Assuming the lemma, the proof of the theorem is concluded as follows.
Since $a(s)$ is cohomologous to $1$, there exists $b\in\GL(N,\hat O')$
such that $b=a(s)s(b)$ for all $s\in G$.  If $b=(b_{ij})$, define a new
basis $e'_1,\ldots,e'_N$ of $W^0_{K'}$ by
\[
e'_j=\sum_{i=1}^{N} b_{ij}e_i.
\]
Using the identity $b=a(s)s(b)$, one sees that $e'_1,\ldots,e'_N$ are
invariant under $G$, hence belong to $W^0$; this proves the
surjectivity of $\hat K'\tens_K W^0 \to W^0_{K'}$.
\end{proof}

\begin{proof}[Proof of the lemma.]
Let $\pi$ be a uniformizing element of $\hat O'$.  Filter the ring
$A=\GL(N,\hat O')$ by means of
$A_n=\{a\in A \mid a\equiv 1 \bmod \pi^n\}$.  We get
%% p125: "GL(N, k'/k)" and "A_n/A_{n+k}" are AS PRINTED (verified at
%% dpi=500); mathematically one expects GL(N,k') and A_n/A_{n+1}.
$A/A_1\simeq\GL(N,k'/k)$, where $k'/k$ is the residue field extension
of $K'/K$.  Moreover, for $n\geq 1$, there is an isomorphism
$A_n/A_{n+k}\simeq M_N(k')$, where $M_N(k')$ is the additive group of
$N\times N$ matrices with coefficients in $k'$.  The lemma follows now
from the triviality of $H^1(G,\GL(N,k'))$ and $H^1(G,M_N(k'))$, where
now $k'$ is endowed with the discrete topology (so this is ordinary
Galois cohomology, cf.~\cite{serre-corps}, p.~158--159).
\end{proof}

\subsechead{A.2}{Admissible characters}

Let $G=\Gal(\bar K/K)$ and let $\phi\colon G\to K^*$ be a continuous
homomorphism.

\begin{dfn}
The character $\phi$ is said to be admissible (notation: $\phi\sim 1$)
if there exists $x\in C$, $x\neq 0$, such that $s(x)=\phi(s)x$ for all
$s\in G$.
\end{dfn}

\begin{remarks}
\begin{nlist}
\item The admissible characters form a subgroup of the group of all
characters of $G$ with values in $K^*$; if $\phi,\phi'$ are two
characters, we write $\phi\sim\phi'$ if $\phi^{-1}\phi'\sim 1$.
\item Let $H^1(G,C^*)$ be the first cohomology group of $G$ with values
in $C^*$ (cohomology being defined by \emph{continuous} cochains, as in
\ref{sec:III.A.1}).  A continuous character $\phi\colon G\to K^*$ is a
$1$-cocycle, hence defines an element $\bar\phi$ of $H^1(G,C^*)$.  One
has $\bar\phi=\bar\phi'$ if and only if $\phi\sim\phi'$.
\item Define a new action of $G$ on $C$ by means of
\[
(s,c)\mapsto \phi(s)s(c)\qquad\qquad s\in G,\ c\in C.
\]
Denote the $C$-$G$-module thus obtained by $C(\phi)$.  Then $\phi$ is
admissible if and only if $C(\phi)$ and $C$ are isomorphic as
$C$-$G$-modules.
\end{nlist}
\end{remarks}

\begin{proposition}[1]\label{prop:III.A.2.1}
Suppose there exists $c\in C^*$ such that $\phi(s)=s(c)/c$ for $s$ in
some open subgroup $N$ of the inertia group of $G$.  Then $\phi$ is
admissible.
\end{proposition}

\begin{proof}
Let $K'/K$ be the subextension of $\bar K/K$ corresponding to $N$; it
is a finite extension of an unramified one.  Let $W=C(\phi)$, as in
Remark 3, and let $W^0$ (resp.\ $W^0_{K'}$) be the subspace of $W$
consisting of elements invariant by $G$ (resp.\ by $N$).  By
hypothesis, $W^0_{K'}$ is $\neq 0$.  Hence, by \ref{sec:III.A.1},
Theorem~\ref{thm:III.A.1.1}, we also have $W^0\neq 0$, and this means
that $\phi$ is admissible, q.e.d.
\end{proof}

Let now $U_C$ be the group of units of $C$, $U^1_C$ the subgroup of
units congruent to $1$ modulo the maximal ideal, and identify
$\bar k^*$ with the group of multiplicative representatives, so that
$U_C=U^1_C\times\bar k^*$, cf.~\cite{serre-corps}, p.~44.  Define the
logarithm map
\[
\log\colon\ U_C \to C
\]
by
\[
\log(x)=0 \qquad\text{if } x\in\bar k^*
\]
\[
\log(x)=\sum_{n=1}^{\infty}(-1)^{n-1}(x-1)^n/n \qquad\text{if }
x\in U^1_C.
\]
This is a continuous homomorphism and even a local isomorphism.

Moreover:

\begin{lemma}
\begin{alist}
\item $\log$ is surjective.
\item The kernel of $\log\colon U_C\to C$ is
$\bar k^*\times\mu_{p^\infty}$, where $\mu_{p^\infty}$ is the set of all
$p^n$-th roots of unity, for $n=1,2,\ldots$
\end{alist}
\end{lemma}

Assertion (a) follows from the fact that $C$ is algebraically closed,
hence that $U_C$ is divisible.

On the other hand, if $u\in U^1_C$ is such that $\log(u)=0$, one has
$\lim\ u^{p^N}=1$, hence, if $N$ is large enough, $u^{p^N}$ belongs to
a subgroup of $U^1_C$ where $\log$ is injective (for instance the
subgroup of elements $x$ with $x\equiv 1 \bmod p^2$).  Hence
$u^{p^N}=1$, and $u\in\mu_{p^\infty}$; this implies (b).

We now apply the log map to the cohomology groups of $G$ with values in
$U_C, C, C^*,\ldots$ (cohomology being, as usual, defined by
\emph{continuous} cochains).  First, since the valuation group of $C$
is $\Q$, we have the exact sequence
\[
1\to U_C \to C^* \to \Q \to 1.
\]
By Tate's theorem (\cite{tate-pdiv}, \S3.3) one has $H^0(G,C^*)=K^*$,
hence an exact sequence
\[
K^* \to \Q \to H^1(G,U_C) \to H^1(G,C^*) \to 0,
\]

%% ------------------------------------------------------------------
%% Printed pages III-37/III-38 follow (absent from this PDF scan;
%% transcribed from working/mcgill_III37.png / mcgill_III38.png).
%% ------------------------------------------------------------------

or, equivalently:
\[
0\to \Q/\Z \xrightarrow{\ \delta\ } H^1(G,U_C)
\xrightarrow{\ i\ } H^1(G,C^*) \to 0.
\]

Let $N=\Ker(\log)$.  We have the exact sequence
\[
H^1(G,N) \xrightarrow{\ j\ } H^1(G,U_C)
\xrightarrow{\ \lambda\ } H^1(G,C),
\]
where $\lambda$ is induced by the log.  Since $H^1(G,C)$ is a
$K$-vector space, the composite $\lambda\circ\delta\colon \Q/\Z \to
H^1(G,C)$ is $0$, hence there is a unique map
\[
L\colon H^1(G,C^*) \to H^1(G,C)
\]
such that $L\circ i=\lambda$.

\begin{proposition}[2]\label{prop:III.A.2.2}
The map $L\colon H^1(G,C^*)\to H^1(G,C)$ is injective.
\end{proposition}

Using the exact sequences above, one sees it is enough to prove that
$i\circ j\colon H^1(G,N)\to H^1(G,C^*)$ is trivial.  But $N$ is a
discrete subgroup of $\bar K^*$, hence $i\circ j$ factors through
$H^1(G,\bar K^*)$, where now $\bar K^*$ is viewed as a discrete group;
by Theorem 90, $H^1(G,\bar K^*)$ is trivial, hence also $i\circ j$,
q.e.d.

Let now $\phi\colon G\to K^*$ be a continuous character.  Since
$\phi(G)$ is compact, it is contained in $U_K$, hence in $U_C$, and
$\log\phi\colon G\to C$ is an additive $1$-cocycle.  Its cohomology
class in $H^1(G,C)$ is $L\bar\phi$, where $\bar\phi$ is the cohomology
class of $\phi$ in $H^1(G,C^*)$.

\begin{proposition}[3]\label{prop:III.A.2.3}
The properties $\phi\sim 1$ and $L\bar\phi=0$ are equivalent.
\end{proposition}

This follows from the injectivity of $L$.

\begin{corollary}
If there exists a non-zero integer $n$ such that $\phi^n\sim 1$, then
$\phi\sim 1$.
\end{corollary}

Indeed, $L\bar\phi = \frac{1}{n}L\bar\phi^{\,n} = 0$.

\runhead{Remark}

Springer has proved that $H^1(G,C)$ is of dimension $1$ over $K$
(cf.~\cite{tate-pdiv}, \S3.3).  Hence, one can take for basis of
$H^1(G,C)$ the element $L\bar\chi$, where $\chi$ is the fundamental
character defined in chap.~\ref{chap:I}, \S\ref{sec:I.1.2}.  In
particular, for any $\phi\colon G\to K^*$, there is an element
$c(\phi)$ of $K$ such that $L\bar\phi=c(\phi)L\bar\chi$; when $K$ is
locally compact, this $c(\phi)$ may be computed explicitly, see
\ref{sec:III.A.6}, Exer.~2.

\subsechead{A.3}{A criterion for local triviality}

From now on, $E$ denotes a subfield of $K$ having the following
properties:
\begin{alist}
\item $E$ contains $\Q_p$ and $[E:\Q_p]<\infty$ (so that $E$ is locally
compact).
\item $K$ contains all $\Q_p$-conjugates of $E$.
\end{alist}
We denote by $\Gamma_E$ the set of all $\Q_p$-embeddings of $E$ in $K$.
Consider a continuous character
\[
\psi\colon\Gal(\bar K/K)\to E^*
\]
with values in $E$.  For each $\sigma\in\Gamma_E$ this gives a
character $\sigma\circ\psi\colon G\to E^*\xrightarrow{\ \sigma\ } K^*$
of $G=\Gal(\bar K/K)$ into $K^*$.

%% NOTE: the book prints "PROPOSITION 3" here although A.2 already has a
%% PROPOSITION 3 (p. III-38).  Numbering as printed; the two are
%% distinguished in cross-references by their section ("Prop. 3 of A2"
%% on p. III-40 vs "Prop. 3 of A3" on p. III-47).
\begin{proposition}[3]\label{prop:III.A.3.3}
The following two properties are equivalent:
\begin{nlist}
\item[\emph{(1)}] $\psi$ is equal to $1$ on an open subgroup of the
inertia group of $G$,
\item[\emph{(2)}] $\sigma\circ\psi\sim 1$ for all $\sigma\in\Gamma_E$.
\end{nlist}
\end{proposition}

\begin{proof}
(1) $\Rightarrow$ (2) is trivial from the result of \ref{sec:III.A.1}
(since we know that admissibility can be seen on an open subgroup of
the inertia group).

(2) $\Rightarrow$ (1).  We use the log map defined in
\ref{sec:III.A.2}.  Note that $\psi$ takes values in the group $U_E$ of
units of $E$, hence $\log\psi\colon G\to E$ is well defined.  Let $I$
be the inertia group of $G$; the subgroup $\log\psi(I)$ of $E$ is
compact, and hence isomorphic to $\Z_p^n$ for some $n$.  If $W$ is the
$\Q_p$-vector subspace of $E$ generated by $\log\psi(I)$, we see that
$\log\psi(I)$ is a lattice in $W$, and $\dim W=n$.  Note that saying
that $\psi$ is equal to $1$ on an open neighbourhood of $1$ in $I$ is
equivalent to saying that $\log\psi(I)=0$ (since
$\log\colon U_E\to E$ is a local isomorphism).  Suppose this is not the
case, i.e.\ suppose that $n\geq 1$.  Choose a $\Q_p$-linear map
$f\colon E\to K$ such that $\dim f(W)=1$; such a map obviously exists.
By Galois theory (independence of characters) the set $\Gamma_E$ is a
basis of $\Hom_{\Q_p}(E,K)$.  Hence, there exist $k_\sigma\in K$ with
\[
f=\sum_{\sigma\in\Gamma_E} k_\sigma\sigma
\]
and we have $f\circ\log\psi = \sum_\sigma k_\sigma\circ\log\psi
= \sum_\sigma k_\sigma\log(\sigma\circ\psi)$.

But by assumption (and Prop.~\ref{prop:III.A.2.3} of
\ref{sec:III.A.2}), the additive $1$-cocycle
$\log(\sigma\circ\psi)\colon G\to K$ is cohomologous to $0$.  Hence the
same is true for $f\circ\log\psi$.  But we may assume (replacing $f$ by
$p^Nf$, with $N$ large, if necessary) that there exists a continuous
homomorphism $F\colon U_E\to U_K$ such that $f\circ\log=\log\circ F$.
We then have $\log(F\circ\psi)=f\circ\log\psi$ and hence
(cf.~Prop.~\ref{prop:III.A.2.3} of \ref{sec:III.A.2}), $F\circ\psi\sim
1$, i.e.\ $F\circ\psi$ is admissible.  But $F\circ\psi$ has now the
property that $F\circ\psi(I)\subset U_K$ is a $p$-adic Lie group of
dimension $1$ (product of $\Z_p$ with a finite group).  This
contradicts a theorem of Tate (\cite{tate-pdiv}, \S3, Th.~2), hence the
result.
\end{proof}

\subsechead{A.4}{The character $\chi_E$}
%% NOTE: the skeleton had "$\xi_E$"; the book prints "The character
%% chi_E" (verified on p. III-40 at dpi=400), hence the correction.

We keep the same hypotheses on $K$ and $E$ as in the previous section.
By class field theory, the group $\Gal(\bar E/E)^{\mathrm{ab}}$ may be
identified with the completion $\hat E^*$ of $E^*$ with respect to the
topology of open subgroups of finite index.  In particular, we have an
exact sequence
\[
1\to U_E \to \Gal(\bar E/E)^{\mathrm{ab}} \to \hat\Z \to 1,
\]
where $\hat\Z\cong\prod_\ell \Z_\ell$ denotes the completion of $\Z$
with respect to the topology of subgroups of finite index (cf.\ for
instance \cite{artin-tate-cft} or \cite{cassels-frohlich}, chap.~VI,
\S2).

Let now $\pi$ be a uniformizing element of $E$.  The image of $\pi$ in
$\Gal(\bar E/E)^{\mathrm{ab}}$ generates a subgroup whose closure is
isomorphic to $\hat\Z$, and this gives an isomorphism:
\[
\Gal(\bar E/E)^{\mathrm{ab}} \simeq U_E\times\hat\Z.
\]
Let $\mathrm{pr}_\pi\colon \Gal(\bar E/E)^{\mathrm{ab}}\to U_E$ be the
projection associated with this decomposition (the Galois extension of
$E$ corresponding to $\Ker(\mathrm{pr}_\pi)$ is the composite of all
finite abelian extensions of $E$ for which $\pi$ is a norm,
cf.~\cite{cassels-frohlich}, p.~144--145).

On the other hand, the inclusion $E\to K$ defines a homomorphism
$\Gal(\bar K/K)\to\Gal(\bar E/E)$, hence also a homomorphism
\[
\mathbf{r}_E\colon G\to\Gal(\bar E/E)^{\mathrm{ab}}.
\]
Define $\chi_{E,\pi}$ (abbr.\ $\chi_E$) to be the composite
homomorphism
\[
G \to \Gal(\bar E/E)^{\mathrm{ab}} \to U_E
\xrightarrow{\ i\ } U_E,
\]
where $i(x)=x^{-1}$ for $x\in U_E$.  Observe that the restriction of
$\chi_E$ to the inertia group of $G$ is $x\mapsto \mathbf{r}_E(x^{-1})$,
and hence is independent of the choice of $\pi$.

\begin{proposition}[4]\label{prop:III.A.4.4}
Let $F_\pi$ be the Lubin--Tate formal group (\cite{lubin-tate}, see
also \cite{cassels-frohlich}, chap.~VI, \S3) associated to $E$ and
$\pi$.  Let $T$ be its Tate-module, which is free of rank $1$ over the
ring $O_E$ of integers of $E$.  The action of $\Gal(\bar K/K)$ on $T$
is given by the character $\chi_E\colon G\to U_E$, defined above.
\end{proposition}

This follows from the main theorem of \cite{lubin-tate} (see also
\cite{cassels-frohlich}, Th.~3, p.~149).

\begin{corollary}
If $E=\Q_p$ and $\pi=p$, then the character $\chi_E$ coincides with the
character $\chi$ defined in chap.~\ref{chap:I}, \S\ref{sec:I.1.2}.
\end{corollary}

Indeed, the Lubin--Tate group is now the multiplicative group $\Gm$ and
its Tate module is the module $T_p(\mu)$ defined in
chap.~\ref{chap:I}, \S\ref{sec:I.1.2}.

\runhead{Remark}

If $K$ is locally compact, we may identify $G^{\mathrm{ab}}$ to
$\hat K^*$ and the character $\chi_E$ is given by
\[
\hat K^* \xrightarrow{\ \mathbf{N}\ } \hat E^*
\xrightarrow{\ \mathrm{pr}_\pi\ } U_E \xrightarrow{\ i\ } U_E,
\]
where $\mathbf{N}=\mathbf{N}_{K/E}$ is the norm map.  [This follows
from the functional properties of the ``reciprocity law'' of local
class field theory.]

In particular, the restriction of $\chi_E$ to the inertia subgroup
$U_K$ of $G^{\mathrm{ab}}$ is $x\mapsto \mathbf{N}_{K/E}(x^{-1})$.

\subsechead{A.5}{Characters associated with Hodge--Tate decompositions}

Retaining the notation of the previous sections, let
$\rho\colon G\to U_E$ be a continuous homomorphism.  Let $V$ be a
one-dimensional vector space over $E$; we make $G$ act on $V$ by
\[
(s,y)\mapsto \rho(s)y,\qquad s\in G,\ y\in V.
\]
Hence $V$ is a $G$-module.  Let $W=C\tens_{\Q_p}V$, where $C=\bar K$ as
before.  This is a $d$-dimensional vector space over $C$, where
$d=[E:\Q_p]$.  Every element $x$ of $E$ defines a $C$-endomorphism
$a_x$ of $W$ by
\[
a_x(\Sigma c_i\tens y_i)=\Sigma c_i\tens xy_i,\qquad
c_i\in C,\ y_i\in V.
\]
We get in this way a representation of $E$ in the $C$-vector space $W$;
note that the action of $a_x$ commutes with the action of $G$.

Let $\sigma\in\Gamma_E$ and put
\[
W_\sigma=\{w \mid w\in W,\ a_x(w)=\sigma(x)w \qquad\text{for all }\
x\in E\}.
\]

\begin{lemma}[1]\label{lem:III.A.5.1}
\begin{alist}
\item Each $W_\sigma$ is a one-dimensional $C$-vector space stable by
$G$.
\item $W$ is the direct sum of the $W_\sigma$'s, $\sigma\in\Gamma_E$.
\item For each $\sigma\in\Gamma_E$, the Galois module $W_\sigma$ is
isomorphic to $C(\sigma\circ\rho)$.
\end{alist}
{\upshape[For the definition of the ``twisted'' module
$C(\sigma\circ\rho)$ see \ref{sec:III.A.2}, Remark 3.]}
\end{lemma}

\begin{proof}
The assertions (a) and (b) are consequences of the well-known fact that
$C\tens_{\Q_p}E$ is a product of $d$ copies of $C$, the projections
$C\tens_{\Q_p}E\to C$ being given by the elements of $\Gamma_E$.

For (c) note that the same decomposition holds for
$V_K=K\tens_{\Q_p}V$, since $K$ contains all the $\Q_p$-conjugates of
$E$; hence for each $\sigma\in\Gamma_E$, there exists a $w\in W_\sigma$
contained in $V_K$.  For such a $w$, say
$w=\Sigma k_i\tens y_i$ ($k_i\in K$, $y_i\in V$) we have
\begin{align*}
s(w)&=\Sigma k_i\tens s(y_i)\\
&=\Sigma k_i\tens \rho(s)y_i\\
&=a_{\rho(s)}w\\
&=\sigma\circ\rho(s)w \qquad\text{since } w \text{ belongs to }
W_\sigma,
\end{align*}
and this implies that $W_\sigma$ is isomorphic to $C(\sigma\circ\rho)$.
\end{proof}

If $\rho_1$ and $\rho_2$ are two characters of $G$ into $K^*$, then we
shall write $\rho_1\equiv\rho_2$ if $\rho_1$ and $\rho_2$ coincide on
an open subgroup of the inertia group of $G$.

\begin{theorem}[2]\label{thm:III.A.5.2}
Let $\rho,V,W$ be as above and, for each $\sigma\in\Gamma_E$, let
$n_\sigma$ be an integer.  The following are equivalent:
\begin{rlist}
\item $\rho\equiv\prod_{\sigma\in\Gamma_E}
\sigma^{-1}\circ\chi_{\sigma E}^{n_\sigma}$
\item $\sigma\circ\rho\sim\chi^{n_\sigma}$ \quad for all
$\sigma\in\Gamma_E$
\item for every $\sigma\in\Gamma_E$ the Galois-module $W_\sigma$ is
isomorphic to $C(\chi^{n_\sigma})$.
\end{rlist}
{\upshape[Recall that $\chi$ is the character defined in
chap.~\ref{chap:I}, \S\ref{sec:I.1.2}, and that $\chi_{\sigma E}$ is
the one attached to the subfield $\sigma E$ of $K$, as in
\ref{sec:III.A.4}.  Note that, since $\chi_{\sigma E}$ restricted to
the inertia group depends only on $\sigma E$, (i) is meaningful.]}
\end{theorem}

\begin{corollary}
$V$ is of Hodge--Tate type if and only if there exist
$n_\sigma\in\Z$ such that
\[
\rho\equiv\prod_{\sigma\in\Gamma_E}
\sigma^{-1}\circ\chi_{\sigma E}^{n_\sigma}.
\]
\end{corollary}

This follows from (iii) and the fact that $W=C\tens V$ is the direct
sum of the $W_\sigma$'s.

\begin{proof}[Proof of Theorem~\ref{thm:III.A.5.2}.]
We prove first:

\begin{lemma}[2]\label{lem:III.A.5.2}
\begin{alist}
\item $\chi_E\sim\chi$
\item If $\sigma\in\Gamma_E$ is not the inclusion map,
$\sigma\chi_E\sim 1$.
\end{alist}
\end{lemma}

\begin{proof}
Let $\pi$ be a uniformizing parameter of $E$, let $F_\pi$ be the
Lubin--Tate group associated to $E$ and $\pi$, let $T_\pi$ be its Tate
module, and $V_\pi=T_\pi\tens\Q_p$.  Since $V_\pi$ is a one-dimensional
vector space over $E$, and $G$ acts on $V_\pi$ through
$\chi_E\colon G\to U_E$ (cf.~\ref{sec:III.A.4},
Prop.~\ref{prop:III.A.4.4}), the above constructions apply to $V_\pi$
and $\chi_E$.  By a theorem of Tate (\cite{tate-pdiv}, \S4, Cor.~2 to
Th.~3), $W_\pi=C\tens_{\Q_p}V_\pi$ has a Hodge--Tate decomposition of
the type
\[
W_\pi=W_\pi(0)\oplus W_\pi(1)
\]
where $\dim W_\pi(0)=d-1$, $\dim W_\pi(1)=1$.  More precisely, Tate
defines canonical isomorphisms $W_\pi(0)=C\tens_K\Hom_E(t',K)$, where
$t'$ is the $(d-1)$-dimensional tangent space of the dual of $F_\pi$,
\[
W_\pi(1)=(C\tens_{\Q_p}V_p(\mu))\tens_K t,
\]
where $t$ is the one-dimensional tangent space to $F_\pi$, and
$V_p(\mu)$ is the $\Q_p$-vector space of dimension $1$ defined in
Chap.~\ref{chap:I}, \S\ref{sec:I.1.2}.  Note that
$C\tens_{\Q_p}V_p(\mu)$
%% p137: the typescript prints "is.isomorphic" (stray period).
is isomorphic to $C(\chi)$, hence one gets an isomorphism
\[
W_\pi(1)\cong C(\chi)\tens_K t.
\]
These isomorphisms commute with the action of $E$.

Since $E$ acts on $t$ by the inclusion map $\sigma_1\colon E\to K$,
this shows that the component $(W_\pi)_{\sigma_1}$ of $W_\pi$ is
$W_\pi(1)$.  Hence, using Lemma~\ref{lem:III.A.5.1}, we have
$C(\chi)\cong C(\chi_E)$, and this implies $\chi_E\sim\chi$, whence
(a).  On the other hand, the same argument shows that
$(W_\pi)_\sigma$, $\sigma\neq\sigma_1$, are contained in the other
factor $W_\pi(0)$ of $W_\pi$; hence $C(\sigma\circ\chi_E)\cong C(1)$,
(where $1$ stands, of course, for the unit character), and this proves
(b).
\end{proof}

We now go back to the proof of Theorem~\ref{thm:III.A.5.2}.  The
equivalence of (ii) and (iii) follows from
Lemma~\ref{lem:III.A.5.1}.  To show (i) $\Longleftrightarrow$ (ii),
note first that $\chi_{\sigma E}$ takes values in $\sigma E^*$, hence
$\sigma^{-1}\circ\chi_{\sigma E}$ takes values in $E^*$, and the same
is true for the character
\[
\rho_1=\prod_{\sigma\in\Gamma_E}
\sigma^{-1}\circ\chi_{\sigma E}^{n_\sigma}.
\]
Let $\tau\in\Gamma_E$.  We have
\[
\tau\circ\rho_1=\prod_{\sigma\in\Gamma_E}
\tau\circ\sigma^{-1}\circ\chi_{\sigma E}^{n_\sigma}.
\]
From Lemma~\ref{lem:III.A.5.2}, applied to the field $\sigma E$, we see
that $\tau\circ\sigma^{-1}\circ\chi_{\sigma E}\sim 1$ if
$\tau\circ\sigma^{-1}$ is not the identity on $\sigma E$, i.e.\ if
$\tau\neq\sigma$; if $\tau=\sigma$, we have
$\tau\circ\sigma^{-1}\circ\chi_{\sigma E}=\chi_{\sigma E}\sim\chi$.
Hence $\tau\circ\rho_1\sim\chi^{n_\tau}$, and (ii) is equivalent to
\[
\tau\circ\rho_1\sim\tau\circ\rho \qquad\text{for all }\
\tau\in\Gamma_E.
\]
By Prop.~\ref{prop:III.A.3.3} of \ref{sec:III.A.3}, this is equivalent
to $\rho_1\equiv\rho$, q.e.d.
\end{proof}

\subsechead{A.6}{Locally compact case}

We now add to all the previous assumptions regarding $K$ and $E$, the
assumption that $K$ is \emph{finite} over $\Q_p$ (i.e.\ $K$ is locally
compact).  By local class field theory, we may then identify
$G^{\mathrm{ab}}$ with $\hat K^*$, and the inertia subgroup of
$G^{\mathrm{ab}}$ with $U_K$, the group of units of $K$.

Let $T$ (resp.\ $T_E, T_{\sigma E}$) be the $\Q_p$-torus associated to
$K$ (resp.\ to $E$, $\sigma E$, where $\sigma\in\Gamma_E$),
cf.~\S\ref{sec:III.1.1}.  The norm map from $K$ to $\sigma E$ defines
an algebraic morphism
\[
\mathbf{N}_{K/\sigma E}\colon T\to T_{\sigma E}.
\]
By composition with $\sigma^{-1}\colon T_{\sigma E}\to T_E$, this gives
a morphism
\[
\mathbf{r}_\sigma=\sigma^{-1}\circ\mathbf{N}_{K/\sigma E}\colon
T\to T_E.
\]

\begin{proposition}[5]\label{prop:III.A.6.5}
\begin{alist}
\item $\mathbf{r}_\sigma(u^{-1})=\sigma^{-1}\circ\chi_{\sigma E}(u)$
for all $u\in U_K$,
\item the $\mathbf{r}_\sigma$ ($\sigma\in\Gamma_E$) make a
$\Z$-basis of $\Hom_{\mathrm{alg}}(T,T_E)$.
\end{alist}
\end{proposition}

(Note that (a) makes sense, since $U_K$ has been identified with the
inertia group of $G^{\mathrm{ab}}$.)

Assertion (a) follows from the remark at the end of \ref{sec:III.A.4}.
On the other hand, let $X(T)$ and $X(T_E)$ be the character groups of
$T$ and $T_E$ respectively.  The characters $[s]$, $s\in\Gamma_K$
(resp.\ $(\sigma)$, $\sigma\in\Gamma_E$) make a basis of $X(T)$
(resp.\ of $X(T_E)$).  The morphism
$\mathbf{r}_\sigma\colon T\to T_E$ defines by transposition a
homomorphism
\[
X(\mathbf{r}_\sigma)\colon X(T_E)\to X(T).
\]
One checks easily that the effect of $X(\mathbf{r}_\sigma)$ on the
basis $[\tau]$, $\tau\in\Gamma_E$, is:
\[
X(\mathbf{r}_\sigma)([\tau])=\sum_{s\sigma=\tau}[s].
\]
Assertion (b) then follows from:

\begin{lemma}
The elements $X(\mathbf{r}_\sigma)$, $\sigma\in\Gamma_E$, form a basis
of $\Hom_{\Gal}(X(T_E),X(T))$.
\end{lemma}

\begin{proof}
The independence of the $X(\mathbf{r}_\sigma)$ is clear.  On the other
hand, let $\phi\in\Hom_{\Gal}(X(T_E),X(T))$ be such that
\[
\phi([\tau])=\Sigma n(\tau,s)[s].
\]
If $\alpha\in\Gal(\bar\Q_p/\Q_p)$ is equal to the identity on
$\tau E$, we have $\alpha[\tau]=[\tau]$, hence
$\alpha\phi([\tau])=\phi([\tau])$, i.e.\ $n(\tau,\alpha s)=n(\tau,s)$
for all $s\in\Gamma_K$.  This means that $n(\tau,s)$ depends only on
the element $\sigma=s^{-1}\tau$; if we put $n_\sigma=n(\tau,s)$, we
then have
\begin{align*}
\phi([\tau])&=\sum_{\sigma\in\Gamma_E} n_\sigma
\sum_{s\sigma=\tau}[s]\\
&=\sum_{\sigma\in\Gamma_E} n_\sigma
X(\mathbf{r}_\sigma)([\tau]).
\end{align*}
This proves the lemma.
\end{proof}

\begin{proposition}[6]\label{prop:III.A.6.6}
Let $\rho$ and $(n_\sigma)$, $\sigma\in\Gamma_E$, be as in
Th.~\ref{thm:III.A.5.2} of \ref{sec:III.A.5}.  Let
$\mathbf{r}\colon T\to T_E$ be the morphism defined by
\[
\mathbf{r}=\prod_{\sigma\in\Gamma_E}
\mathbf{r}_\sigma^{n_\sigma}.
\]
The equivalent properties (i), (ii), (iii) of
Th.~\ref{thm:III.A.5.2} are equivalent to:
\begin{nlist}
\item[\emph{(iv)}] There exists an open subgroup $U'$ of the inertia
subgroup $U_K$ of $G^{\mathrm{ab}}$ such that
$\mathbf{r}(u)\rho(u)=1$ if $u\in U'$.
\end{nlist}
\end{proposition}

Indeed, (iv) is just a reformulation of (i), since we know that
$\sigma^{-1}\circ\chi_{\sigma E}(u)=\mathbf{r}_\sigma(u^{-1})$ if
$u\in U_K$.

\begin{corollary}
The following are equivalent:
\begin{alist}
\item $\rho$ is locally algebraic.
\item The Galois module $V$ attached to $\rho$ is of Hodge--Tate type.
\end{alist}
\end{corollary}

This follows from Theorem~\ref{thm:III.A.5.2}, combined with
Prop.~\ref{prop:III.A.6.5} and Prop.~\ref{prop:III.A.6.6}.

\begin{exercises}

\begin{nlist}

\item a) Let $A=\End_{\Q_p}(K)$ be the space of $\Q_p$-linear
endomorphisms of $K$; if $a\in A$, denote by $\mathrm{Tr}(a)$ the trace of
$a$.  If $x\in K$, denote by $u_x$ the endomorphism $y\mapsto xy$ of
$K$.  Show that, for any $a\in A$, there exists a unique element
$c_K(a)$ of $K$ such that
\[
\mathrm{Tr}(u_x\circ a)=\mathrm{Tr}_{K/\Q_p}(x\cdot c_K(a))
\qquad\text{for all } x\in K.
\]

b) Show that the map $c_K\colon A\to K$ so defined is $K$-linear for
both the natural structures of $K$-vector space on $A$.

c) Let $e_i$ be a $\Q_p$-basis of $K$ and let $e'_i$ be the dual
basis, so that $\mathrm{Tr}_{K/\Q_p}(e_i e'_j)=\delta_{ij}$.  Show that
\[
c_K(a)=\sum_i a(e_i)e'_i \qquad\text{if } a\in A.
\]

d) If $L\supset K$ and $a\in A$, show that
\[
c_L(a\circ\mathrm{Tr}_{L/K})=c_K(a).
\]
Show that $c_K(\mathrm{Tr}_{K/\Q_p})=1$.

e) If $K$ is a Galois extension of $\Q_p$, show that $c_K(\sigma)=0$
for every $\sigma\in\Gal(K/\Q_p)$, $\sigma\neq$ id., and
$c_K(\mathrm{id.})=1$.

\item Let $\phi\colon G^{\mathrm{ab}}\to K^*$ be a continuous
homomorphism, and let $a_\phi$ be the $\Q_p$-linear endomorphism of
$K$ such that the diagram
\[
\begin{array}{ccc}
U_K & \xrightarrow{\ \phi\ } & U_K\\[0.3ex]
\downarrow\log && \downarrow\log\\[0.3ex]
K & \xrightarrow{\ a_\phi\ } & K
\end{array}
\]
is commutative.  Let $L\bar\phi$ (resp.\ $L\bar\chi$) be the image of
$\phi$ (resp.\ $\chi$) in the one-dimensional $K$-vector space
$H^1(G,C)$, cf.~\ref{sec:III.A.2}.  Show that
\[
L\bar\phi=c\cdot L\bar\chi
\]
where $c=-c_K(a_\phi)$.  (Check the formula first when $K$ is a Galois
%% p143: "Gal(K(Q_p))" is AS PRINTED; evidently for Gal(K/Q_p).
extension of $\Q_p$ and $\phi=\sigma^{-1}\circ\chi_K$,
$\sigma\in\Gal(K(\Q_p))$, in which case $a_\phi=-\sigma^{-1}$ and
$c_K(a_\phi)$ is given by Exer.~1, d.)

In particular, $\phi$ is admissible if and only if $c_K(a_\phi)=0$.

\end{nlist}
\end{exercises}

\subsechead{A.7}{Tate's theorem}

We recall the statement (cf.~\S\ref{sec:III.1.2}); here again, $K$ is
locally compact.

\begin{theorem}[3]\label{thm:III.A.7.3}
Let $V$ be a finite dimensional vector space over $\Q_p$ and let
%% p143: the superscript "al" in "G^{al}" is AS PRINTED (verified at
%% dpi=400); evidently for "G^{ab}", since rho is an abelian
%% representation.
$\rho\colon G^{\mathrm{al}}\to\Aut(V)$ be an abelian $p$-adic
representation of $K$.  The following are equivalent:
\begin{nlist}
\item[\emph{(1)}] $\rho$ is locally algebraic
\item[\emph{(2)}] $\rho$ is of Hodge--Tate type and its restriction to
the inertia group is semi-simple.
\end{nlist}
\end{theorem}

\begin{proof}
We have already remarked (cf.~\S\ref{sec:III.1.1}) that (1) implies:

{\upshape(*)\ --- \emph{The restriction of $\rho$ to the inertia group
is semi-simple.}}  Hence we may assume that ($*$) holds.

Let $\pi$ be a uniformizing element of $K$, and let
$\mathrm{pr}_\pi$ denote the projection map of $G^{\mathrm{ab}}$ onto
its inertia group $U_K$ associated to $\pi$ (cf.~\ref{sec:III.A.4} and
\cite{cassels-frohlich}, p.~144--145).  Define a new representation
$\rho'$ of $G^{\mathrm{ab}}$ by
\[
\rho'=\rho\circ\mathrm{pr}_\pi.
\]
Replacing $\rho$ by $\rho'$ does not affect the local algebraicity
(clear), nor the Hodge--Tate property (this follows from
\ref{sec:III.A.1}, Cor.~\ref{cor:III.A.1.2} to
Th.~\ref{thm:III.A.1.1}).  Since ($*$) implies that $\rho'$ is
semi-simple, this means that, after replacing $\rho$ by $\rho'$, we may
assume that $\rho$ is semi-simple and even (by an easy reduction) that
it is \emph{simple}.  Let then $E\subset\End(V)$ be the commuting
algebra of $\rho$.  Since $\rho$ is abelian and simple, $E$ is a
commutative field, of finite degree over $\Q_p$, and $V$ is a
one-dimensional vector space over $E$; the representation $\rho$ is
given by a continuous character $\rho\colon G\to E^*$.

Let now $K'$ be a finite extension of $K$ which is large enough to
contain all the $\Q_p$-conjugates of $E$.  Call (1$'$) and (2$'$) the
properties corresponding to (1) and (2), when $K'$ is taken as
groundfield instead of $K$.  We know (cf.~\S\ref{sec:III.1.1}) that
(1) $\Longleftrightarrow$ (1$'$).  By
Cor.~\ref{cor:III.A.1.2} to Th.~\ref{thm:III.A.1.1} of
\ref{sec:III.A.1}, we have (2) $\Longleftrightarrow$ (2$'$).  Hence it
is enough to prove that (1$'$) $\Longleftrightarrow$ (2$'$), and this
has been done in \ref{sec:III.A.6} (Cor.~to
Prop.~\ref{prop:III.A.6.6}), q.e.d.
\end{proof}
